\lecture{Lecture 8: Detecting edges and lines}

\qa{How do first-order and second-order derivatives differ in their approach to edge detection?}{
Derivatives are used to detect rapid spatial changes in pixel intensity, which define edges.
First-order derivatives (like the Sobel operator) compute the gradient of the image. An edge is identified by locating the local maxima (the highest peaks) in the gradient magnitude. The gradient also provides the direction of the edge.
Second-order derivatives (like the Laplacian) compute the rate of change of the gradient. An edge is identified by locating zero-crossings—the exact point where the second derivative passes through zero. While second-order derivatives are better at pinpointing the exact center of an edge, they are exponentially more sensitive to high-frequency noise. Both derivatives amplify noise. For this reason, it's strictly necessary to smooth the image before applying edge detection. The most appropriate and used choice for this is a Gaussian Low-Pass Filter, which has a smooth derivative and it's also the foundation of Canny Edge Detector.
}

\qa{In the Canny algorithm pipeline, what distinct structural problems do Non-Maxima Suppression and Hysteresis Thresholding solve?}
{\begin{itemize}
    \item \textbf{Non-Maxima Suppression} solves the problem of thick, blurred edge responses. It ensures precise edge localization by thinning the gradient down to a single-pixel line. It achieves this by suppressing (setting to 0) any pixel that is not the local maximum along its quantized gradient direction.
    \item \textbf{Hysteresis Thresholding} solves the problem of fragmented edges and noise. By using two thresholds ($T_H$ and $T_L$), it keeps strong edges, completely rejects isolated weak responses (noise), and most importantly, preserves weak edge pixels \textit{only} if they are spatially connected to a strong edge, thus closing gaps.
\end{itemize}}

\qa{How does the choice of the Gaussian kernel size (\texorpdfstring{$\sigma$}{sigma}) in Step 1 dictate the final output of the Canny edge detector?}
{\begin{itemize}
    \item \textbf{A small $\sigma$} applies minimal smoothing. This allows the detector to capture finer details and highly localized edges, though it makes the output more susceptible to keeping small noise components.
    \item \textbf{A large $\sigma$} applies aggressive smoothing. This blurs out fine textures and noise, restricting the detector to finding only coarse, large-scale, and highly prominent structural edges in the image.
\end{itemize}}

\qa{Why is Canny considered a complete "edge detector algorithm" rather than just a simple derivative filter?}
{\begin{itemize}
    \item A \textbf{simple derivative filter} (e.g., computing the gradient magnitude and applying a basic threshold) merely highlights areas of high intensity change. As seen in the slides, this typically results in thick, noisy, and broken boundary responses.
    \item The \textbf{Canny edge detector} embeds a complex, multi-step mathematical model of what an edge should be. It actively enforces three strict design criteria: a low error rate (via Gaussian smoothing), precise single-point localization (via non-maxima suppression), and continuous contour extraction (via hysteresis thresholding).
\end{itemize}}

\qa{When building the parameter space for the Hough Transform, why is the normal representation (\texorpdfstring{$x \cos\theta + y \sin\theta = \rho$}{x cos(theta) + y sin(theta) = rho}) preferred over the standard Cartesian equation (\texorpdfstring{$y = ax + b$}{y=ax+b})?}
{\begin{itemize}
    \item \textbf{Cartesian Representation}: Maps an image point $(x, y)$ to a line $b = -xa + y$ in the $ab$-parameter space. However, it fails to represent vertical lines because their slope ($a$) approaches infinity.
    \item \textbf{Normal Representation}: Solves the vertical line issue by mapping an image point to a bounded sinusoidal curve in the $\rho\theta$-parameter space. This allows the algorithm to gracefully detect lines at any angle.
\end{itemize}}

\qa{In the Hough Transform algorithm, the parameter space (\texorpdfstring{$\rho$}{rho}, \texorpdfstring{$\theta$}{theta}) must be quantized into accumulation cells. What is the engineering trade-off between using few cells versus many cells?}
{\begin{itemize}
    \item \textbf{Few Cells (Coarse Quantization)}: Makes the algorithm robust to noise, allowing it to handle pixels that are not perfectly aligned along a mathematical line. The drawback is poor line localization.
    \item \textbf{Many Cells (Fine Quantization)}: Yields very accurate line localization, but requires the edge pixels in the image to be precisely aligned. If they are slightly scattered, they won't accumulate in the same cell, and the line might be missed.
\end{itemize}}

\qa{How does the Generalized Hough Transform differ conceptually from the standard version, and what is its main computational drawback?}
{\begin{itemize}
    \item \textbf{Standard Hough Transform}: Detects straight lines by finding intersections in a 2D parameter space (either $ab$ or $\rho\theta$).
    \item \textbf{Generalized Hough Transform}: Extends the concept to find more complex shapes (like circles) using a general equation $g(\boldsymbol{v}, \boldsymbol{c}) = 0$, where $\boldsymbol{v}$ are coordinates and $\boldsymbol{c}$ are coefficients. For example, a circle uses $(x - c_1)^2 + (y - c_2)^2 = c_3^2$. The main drawback is that complex shapes require a parameter space with \textbf{high dimensionality} (e.g., 3 dimensions for a circle), drastically increasing the computational and memory cost.
\end{itemize}}

\qa{Why do we use the Hough Transform instead of a naive approach to find lines? (Computational Complexity)}
{\begin{itemize}
    \item A \textbf{naive approach} would consider all couples of edge points, evaluate the line passing through them, and count the points on that line. This yields a heavy computational complexity of $O(n^3)$ (considering $O(n^2)$ couples plus the comparisons).
    \item The \textbf{Hough Transform} avoids this combinatorics explosion by mapping each edge pixel independently into a parameter space and simply voting (accumulating) in a discrete grid, making line detection highly efficient.
\end{itemize}}